% ----------------------------------------------------------
% Introdução
% ----------------------------------------------------------
\chapter{Introdução}

Este template é apenas um exemplo para auxiliar na elaboração do trabalho, não descarta a consulta no Guia de trabalhos acadêmicos da instituição.

A introdução abre o trabalho propriamente dito. Tem a finalidade de apresentar os motivos que levaram o autor a realizar a pesquisa, o problema abordado, os objetivos e a justificativa. O objetivo principal da introdução é situar o leitor no contexto da pesquisa. O leitor deverá perceber claramente o que foi analisado, como e por que, as limitações encontradas, o alcance da investigação e suas bases teóricas gerais. Ela tem, acima de tudo, um caráter didático de apresentar o que foi investigado, levando-se em conta o leitor a que se destina e a finalidade do trabalho.

Assim, na introdução contextualize o tema, delimite o assunto, apresente um rápido histórico do problema e das soluções porventura já apresentadas, com breve revisão crítica das investigações anteriores; faça referência às fontes de material, aos métodos seguidos, às teorias ou aos conceitos que embasam o desenvolvimento e a argumentação, às eventuais faltas de informação, ao instrumental utilizado.

A introdução deverá conter, ainda: Justificativa; Definição do problema; Objetivo geral e objetivos específicos.

Todos os autores citados devem ter a referência incluída em lista no final no trabalho. \cite{de2011associaccao}

\section{Justificativa}
%espaço para a justificativa do trabalho

\section{Definição do problema}
%espaço para definição do problema do trabalho

\section{Objetivo geral}
%espaço para os objetivos do trabalho

\subsection{Objetivos específicos}